\section{Theoretical guarantees}
\label{sec:theory}

This section states the three results used by FORGE: a forward-forgetting rate on the
feature block (Theorem~\ref{thm:forward-forgetting}), a subtraction bound on the head
(Lemma~\ref{lem:subtraction}), and a composition theorem for the combined protocol
(Theorem~\ref{thm:composition}). Proofs are in Appendix~\ref{app:proofs}.

\begin{theorem}[Forward-forgetting rate]
\label{thm:forward-forgetting}
Consider Algorithm~\ref{alg:decayed-aggregation} with decay $\alpha \in (0,1)$, clip
$C$, per-round feature noise $\sigma$, and server step $\eta_g$. Let client $\forget$
participate in rounds $S \subseteq [0, r_\forget]$ with $|S| \le \rho \, r_\forget$.
Then at round $r_\mathrm{req} = r_\forget + T$, equivalently $\mu_1(T)$-GDP (see
Theorem~\ref{thm:composition}),
\begin{align}
\theta^{(r_\mathrm{req})}_\mathrm{sh}
 &\;\indist[(\eps_1(T), \delta)]\;
\theta^{(r_\mathrm{req})}_\mathrm{sh,-\forget}, \nonumber \\
\eps_1(T) &\;=\;
\frac{C\,\alpha^T\,\sqrt{2\ln(1.25/\delta)}}{\sigma \sqrt{r_\mathrm{req}}}.
\label{eq:eps-rate}
\end{align}
\end{theorem}

The bound combines the $\alpha^T$ decay of a departed client's residual contribution to
the EMA aggregation state with a Gaussian-mechanism release of the trajectory under
accumulated noise variance $r_\mathrm{req}\,\sigma^2$. Equivalently, the feature-block
release is a Gaussian mechanism with sensitivity $C\alpha^T$ and noise scale
$\sigma \sqrt{r_\mathrm{req}}$, which satisfies $\mu_1(T)$-GDP
\cite{dong2019gaussian} with $\mu_1(T) = C \alpha^T / (\sigma \sqrt{r_\mathrm{req}})$.
At the paper's hyperparameters ($\alpha{=}0.5$, $C{=}50$, $\sigma{=}10^{-3}$,
$\delta{=}10^{-5}$), the bound is vacuous for $T \lesssim 20$ and tight for
$T \gtrsim T^\star \approx 50$; a full table is in Appendix~\ref{app:numerical}. The
sensitivity argument is linear in the aggregation state; a trajectory-level Lipschitz
treatment of the non-linear local-training step follows \cite{neel2021descent,
gupta2021adaptive}. The empirical per-round Lipschitz constant $L'$ of the
local-training map, measured directly at the operating point
(Appendix~\ref{app:nonconvex}), has p95 $87$ and max $88$ across 32 (client,
direction) pairs, which satisfies the contraction condition
$\eta_g (1-\alpha) \eta_l L' < 1$ required by the trajectory argument and inflates
the sensitivity by at most a factor of $\approx 9$. The $\alpha^T$ decay dominates
this correction for any $T \gtrsim 20$, so the operating-point numerics in
Theorem~\ref{thm:composition} are unaffected to within this constant.

\begin{lemma}[Head-subtraction cost]
\label{lem:subtraction}
Let $\sigma_\forget$ be the pre-committed summary of client $\forget$'s data, released
by the Gaussian mechanism at per-release $(\eps_\summary, \delta_\summary)$-DP. Let
$\mathcal{E}: (\theta, \sigma) \mapsto \theta'$ be the erasure map of
Algorithm~\ref{alg:erasure} restricted to the head. Then
$\theta'_\mathrm{cl} \indist[(\eps_\summary, \delta_\summary)]
\theta'_\mathrm{cl, -\forget}$.
\end{lemma}

The map $\mathcal{E}(\theta, \cdot)$ is a deterministic post-processing of
$\sigma_\forget$; the bound follows from DP post-processing \cite{dwork2014book}. The
conditional independence of $\theta_\mathrm{cl}$ from $\sigma_\forget$ during training
holds by construction of Algorithm~\ref{alg:decayed-aggregation}: the head is updated
by unweighted FedAvg, which never touches the summary.

\begin{theorem}[Composition]
\label{thm:composition}
Under Algorithm~\ref{alg:decayed-aggregation}, the feature-block release is
$\mu_1(T)$-GDP with $\mu_1(T) = C \alpha^T / (\sigma \sqrt{r_\mathrm{req}})$, and,
conditional on the feature-block release, the head-subtraction release is
$\mu_\summary$-GDP with $\mu_\summary = \Delta / \sigma_\summary$. By adaptive
Gaussian-DP composition \cite{dong2019gaussian}, the joint release
$\theta^{(r_\mathrm{req})}$ satisfies $\mu_\mathrm{joint}(T)$-GDP with
\[
\mu_\mathrm{joint}(T) \;=\; \sqrt{\mu_1(T)^2 + \mu_\summary^2},
\]
which converts to $(\eps_\mathrm{total}, \delta_\mathrm{total})$-DP via the standard
GDP $\leftrightarrow (\eps, \delta)$ duality.
\end{theorem}

\begin{proof}[Proof]
Let $M_1$ denote the feature-block mechanism of
Algorithm~\ref{alg:decayed-aggregation} that releases
$\theta_\mathrm{sh}^{(r_\mathrm{req})}$, and let $M_2$ denote the head-subtraction
mechanism of Algorithm~\ref{alg:erasure} that releases
$\theta_\mathrm{cl}^{(r_\mathrm{req})\prime}$. By
Theorem~\ref{thm:forward-forgetting}, $M_1$ is a Gaussian mechanism with sensitivity
$C\alpha^T$ and noise scale $\sigma \sqrt{r_\mathrm{req}}$, hence $\mu_1(T)$-GDP.
Conditional on the output of $M_1$, the summary release used by $M_2$ is a Gaussian
mechanism on $D_\forget$ with sensitivity $\Delta$ and noise scale $\sigma_\summary$
(Equation~\eqref{eq:summary-moments}); by Lemma~\ref{lem:subtraction} and the
post-processing property of GDP, $M_2$ is $\mu_\summary$-GDP. The feature-block server
noise $\xi^{(t)}$ driving $M_1$ (Algorithm~\ref{alg:decayed-aggregation} line~7) and
the summary noise $\eta^c$ driving $M_2$ (Equation~\eqref{eq:summary-moments}) are
drawn from independent random sources by construction, which is the independence
condition that adaptive GDP composition requires; disjointness of input data is not.
The joint mechanism $(M_1(D_\forget), M_2(D_\forget, M_1(D_\forget)))$ is therefore
the adaptive composition of two Gaussian mechanisms on the same dataset. By
Theorem~3.2 of \cite{dong2019gaussian}, the composed mechanism is
$\mu_\mathrm{joint}(T)$-GDP with $\mu_\mathrm{joint}(T) = \sqrt{\mu_1(T)^2 +
\mu_\summary^2}$. The adaptive clause handles the dependence of $\sigma_\forget$ on
$\theta_\mathrm{sh}$ through the feature embedding $\phi(x_j)$: $\mu_\summary$ is
computed under the feature-block state released by $M_1$, and the composition result
holds pointwise in that state. The released model $\theta^{(r_\mathrm{req})}$ is a
deterministic post-processing of this adaptive composition, which preserves the GDP
parameter; the GDP $\leftrightarrow (\varepsilon, \delta)$ duality yields the stated
$(\varepsilon_\mathrm{total}, \delta_\mathrm{total})$-DP conversion.
\end{proof}

At the paper's hyperparameters with $r_\mathrm{req} - r_\forget = 200$,
$\mu_1(200) \approx 10^{-28}$ (zero at machine precision), so
$\mu_\mathrm{joint}(200) \approx \mu_\summary \approx 0.413$ and the overall guarantee
is $(\eps_\mathrm{total}, \delta_\mathrm{total}) = (2.0, 10^{-5})$.

Two observations on this operating point are in order. First, the role of the decayed
aggregation in the composition is not to contribute numerically to the headline bound
at $T = 200$ but to fix the asymptote: Theorem~\ref{thm:forward-forgetting} drives
$\mu_1(T) \to 0$ as $T \to \infty$ at an exponential rate in $\alpha^T$, which ensures
that the overall guarantee $\sqrt{\mu_1(T)^2 + \mu_\summary^2}$ does not degrade with
the cooling-off period. Without the decay, a static client contribution would appear in
the final-model release, and the overall bound would then grow with $T$ via the
accumulated feature-block exposure. Second, the regime in which the decay contributes
materially to the formal bound is characterised by a short cooling-off period
($T \lesssim T^\star$) under tight hyperparameters, i.e., $\alpha$ close to 1 or
$\sigma$ at the small end of the usable range. The $\alpha$-sweep in
Section~\ref{sec:results-ablation} reports this regime empirically; outside of it, the
headline numerical guarantee is supplied entirely by the summary-subtraction cost
$\mu_\summary$.

The storage model of Section~\ref{sec:method-summary} (hashes plus one most-recent
summary, with purge on new release) ensures that only a single DP release of
$\sigma_\forget$ is observed at erasure time. The summary-channel DP cost is therefore
$\eps_\summary$ per erasure event, and does not compose across the $\lfloor
r_\mathrm{req} / F \rfloor$ release points during training.
